\documentclass[9pt]{extarticle}
\usepackage[T1]{fontenc}
\usepackage{lmodern,amsmath,amssymb,amsthm,microtype}
\usepackage[margin=0.76in]{geometry}
\usepackage[colorlinks=true,urlcolor=blue,citecolor=blue,linkcolor=blue]{hyperref}
\newtheorem{theorem}{Theorem}
\newtheorem{lemma}{Lemma}
\newtheorem{corollary}{Corollary}
\newcommand{\cat}{\mathbin{\Vert}}
\title{Arbitrarily Long Increasing Chains of Squares Whose Concatenation Remains a Square}
\author{Qichao Wang\\{\small Hebei University of Technology, Tianjin, China}\\[0.4em]Daoyu Dong\\{\small University of Electronic Science and Technology of China, Chengdu, China}}
\date{October 17, 2026}\linespread{0.98}\setlength{\abovedisplayskip}{6pt}\setlength{\belowdisplayskip}{6pt}
\begin{document}
\maketitle
\begin{abstract}
We give an elementary, constructive existence theorem for concatenations of
perfect squares. Every positive square $X^2$ can be extended to a larger
square by appending a nonzero square having an arbitrarily large prescribed
even digit length from an infinite set of lengths. The appended block may
be arbitrarily large. Iteration proves that for each $n\ge1$ there exist
strictly increasing positive squares $s_1<\cdots<s_n$ whose decimal
concatenation $s_1\cat\cdots\cat s_n$ is a square. In particular,
OEIS A087354, A391022, and A391189 are defined at every positive index.
The argument uses factorization of a difference of squares and the density
of the multiplicative group $\{2^u5^v:u,v\in\mathbb Z\}$ in the
positive reals. It does not determine the minimum values of the sequences.
\end{abstract}

\section{Background and statement}
For positive integers $a_1,\ldots,a_n$, let
$a_1\cat\cdots\cat a_n$ denote their decimal concatenation with no
leading zeros in any block. OEIS \href{https://oeis.org/A087354}{A087354}
records the \emph{least} square representable as a concatenation of
$n$ nonzero squares; repetitions are permitted \cite{oeis87354}.
The more restrictive \href{https://oeis.org/A391022}{A391022}
requires the blocks to be distinct and strictly increasing, while
\href{https://oeis.org/A391189}{A391189} minimizes the largest
square root in such a representation \cite{oeis391022,oeis391189}.
The issue addressed here is existence for \emph{every} $n$, not
optimization of any of these minima.

\begin{theorem}[Unbounded extension of an arbitrary square]\label{thm:extend}
For each positive integer $X$ and each real bound $M>0$, there exist
positive integers $h,Y,Z$ such that
\begin{equation}\label{eq:extend}
 Z^2=X^2 10^{2h}+Y^2,\qquad
 10^{2h-1}<Y^2<10^{2h},\qquad Y^2>M.
\end{equation}
In particular, appending the decimal representation of the nonzero
square $Y^2$ to that of $X^2$ yields exactly the square $Z^2$.
For each fixed $X$, there are infinitely many such extensions.
\end{theorem}

\begin{corollary}[All-index existence]\label{cor:oeis}
For every integer $n\ge1$, there exist positive integers
\[s_1<s_2<\cdots<s_n\]
such that every $s_i$ is a square and
$s_1\cat\cdots\cat s_n$ is a square. Consequently, each of
A087354$(n)$, A391022$(n)$ and A391189$(n)$ exists and is positive.
There are infinitely many such concatenations for each $n\ge2$.
\end{corollary}
\begin{proof}
Start with $s_1=1$ and $X_1=1$, so $X_1^2=s_1$.
Suppose $s_1<\cdots<s_j$ have been chosen and
$X_j^2=s_1\cat\cdots\cat s_j$. Apply
Theorem~\ref{thm:extend} to $X=X_j$ with $M=s_j$.
Set $s_{j+1}=Y^2$ and $X_{j+1}=Z$. The digit-length
inequality in \eqref{eq:extend} ensures that this is precisely
concatenation (not addition with spurious leading zeros), and
$s_{j+1}>s_j$. Induction proves the claim.
At every extension step one can choose arbitrarily large $h$,
producing infinitely many distinct concatenations. Since the set of
eligible concatenated squares for each $n$ is a nonempty set of
positive integers, its least member exists; similarly, the set
of eligible largest square roots has a least member.
\end{proof}

\section{The construction}
\begin{lemma}[Density of decimal-smooth rational ratios]\label{lem:dense}
The multiplicative subgroup
\[\mathcal G=\{2^u5^v:u,v\in\mathbb Z\}\]
is dense in $(0,\infty)$.
\end{lemma}
\begin{proof}
The ratio $\log 2/\log 5$ is irrational: if it were $a/b$
for nonzero integers $a,b$, then $2^b=5^a$, contradicting
unique factorization (with signs dealt with by inversion).
The fractional parts of integer multiples of an irrational number
are dense in $[0,1]$. Thus
$\{u\log2+v\log5:u,v\in\mathbb Z\}$
is dense in $\mathbb R$: for a target $t$, approximate
$u(\log2/\log5)$ modulo 1 to $t/\log5$ modulo 1,
and then choose $v$ to adjust the integer part.
Exponentiation establishes the assertion.
\end{proof}

\begin{proof}[Proof of Theorem~\ref{thm:extend}]
Fix $X\ge1$. For $q>1$ put
\[f_X(q)=\frac X2(q-q^{-1}).\]
This is a continuous strictly increasing function from $0$ to
$+\infty$ as $q$ increases from $1$ to $+\infty$. Therefore
\begin{equation}\label{eq:openinterval}
\mathcal I_X=\left\{q>1:\frac1{\sqrt{10}}<f_X(q)<1\right\}
\end{equation}
is a nonempty open interval. By Lemma~\ref{lem:dense}, it
contains a rational number of the special form $q=2^u5^v$
with $u,v\in\mathbb Z$.

Fix this $q$, and take any integer
$h\ge\max\{|u|,|v|\}+1$. Define two positive even integers
\begin{equation}\label{eq:factors}
 A=X10^h q^{-1}=X2^{h-u}5^{h-v},\qquad
 B=X10^h q=X2^{h+u}5^{h+v}.
\end{equation}
The exponents on the right are nonnegative integers and, because
$h\ge |u|+1$, both $A$ and $B$ are even. Since $q>1$, $B>A$.
Consequently
\begin{equation}\label{eq:yz}
 Y=\frac{B-A}{2}\in\mathbb Z_{>0},\qquad
 Z=\frac{B+A}{2}\in\mathbb Z_{>0}
\end{equation}
are integers, and
\begin{equation}\label{eq:pyth}
 Z^2-Y^2=AB=X^2 10^{2h}.
\end{equation}
Moreover, \eqref{eq:openinterval} gives
\[\frac{Y}{10^h}=f_X(q)\in\left(\frac1{\sqrt{10}},1\right).\]
Squaring,
\[10^{2h-1}<Y^2<10^{2h},\]
so $Y^2$ has exactly $2h$ decimal digits and is nonzero.
The same fixed $q$ works for all larger $h$, and
$Y^2=10^{2h}f_X(q)^2\to\infty$. Taking $h$ arbitrarily
large forces $Y^2>M$ and yields infinitely many extensions.
\end{proof}

\section{Examples, variants, and scope}
For $X=1$, the choice $q=2$, $h=2$ gives
\[
 A=50,\quad B=200,\quad Y=75,\quad Z=125,
 \qquad\boxed{1\cat5625=15625=125^2}.
\]
An example starting from the strictly increasing two-block
square $4\cat9=49=7^2$ is obtained using
$X=7$, $q=2^{14}5^{-6}=16384/15625$, and $h=15$:
\[
Y=332155892578125,\qquad Z=7007876107421875,
\]
\begin{equation}\label{eq:example}
\boxed{4\cat9\cat110327536974370914459228515625
=7007876107421875^2.}
\end{equation}
The third block is the $30$-digit square of $Y$ and is
larger than $9$. All stated equalities can be checked with
integer arithmetic. The accompanying script uses rational rather
than floating-point inequalities to verify the construction.

The proof works in any integer base $b\ge2$ that has at least
two distinct prime factors. Indeed, choose two such primes
$\rho,\sigma\mid b$ and replace $q$ by $\rho^u\sigma^v$, $10^h$
by $b^h$, and $1/\sqrt{10}$ by $1/\sqrt b$. The subgroup generated by
$\rho,\sigma$ is dense for the same unique-factorization reason;
for sufficiently large $h$ the factors $Xb^hq$ and $Xb^h/q$
are integral and of the same parity (both odd if $bX$ is odd;
both even otherwise). Thus the same extension and induction work
in those bases. We do not claim this argument for prime-power bases.

These are \emph{existence} theorems and may produce very large
witnesses. They do not provide the exact minimal
values in the three OEIS entries. The related A391022 and A391189 entries
list only their early values; neither entry currently states
an explicit universal existence theorem. Literature priority
beyond the references below remains to be checked independently.

\begin{thebibliography}{9}
\bibitem{oeis87354} OEIS Foundation Inc.,
\emph{A087354: Smallest square which is the concatenation of $n$
nonzero squares}, \url{https://oeis.org/A087354}, accessed October 17, 2026.
\bibitem{oeis391022} OEIS Foundation Inc.,
\emph{A391022: Minimal square obtained by concatenating $n$ distinct
nonzero squares in ascending order}, \url{https://oeis.org/A391022},
accessed October 17, 2026.
\bibitem{oeis391189} OEIS Foundation Inc.,
\emph{A391189: Least upper bound on square roots of $n$ distinct
ascending square blocks whose concatenation is square},
\url{https://oeis.org/A391189}, accessed October 17, 2026.
\end{thebibliography}
\end{document}
